%================================================================
% MAA MathFest 2026 Undergraduate Poster -- 36 in x 24 in
% Christian Dane Beels, advised by Joseph Dorta, Ph.D.
%================================================================

\documentclass[25pt,a0paper,portrait]{tikzposter}
\geometry{paperwidth=36in,paperheight=24in,margin=0mm}

% tikzposter computes these lengths before the custom page geometry above.
% Recompute them so the column grid matches the actual 36 x 24 inch sheet.
\makeatletter
\setlength{\TP@innermargin}{8mm}
\setlength{\TP@colspace}{8mm}
\setlength{\TP@blockverticalspace}{8mm}
\setlength{\TP@visibletextwidth}{\textwidth-2\TP@innermargin}
\setlength{\TP@visibletextheight}{\textheight-2\TP@innermargin}
\makeatother

\usepackage[utf8]{inputenc}
\usepackage[T1]{fontenc}
\usepackage{amsmath,amssymb}
\usepackage{enumitem}
\usepackage{relsize}

\setlist[itemize]{leftmargin=*,itemsep=4pt,topsep=3pt,parsep=0pt,partopsep=0pt}

\newcommand{\bodyfont}{\fontsize{25.5}{30.5}\selectfont}
\newcommand{\smallbody}{\fontsize{25}{30.5}\selectfont}
\newcommand{\captionfont}{\fontsize{21.5}{25.5}\selectfont}
\newcommand{\referencefont}{\fontsize{18}{21.5}\selectfont}
\newcommand{\theorybody}{\fontsize{26}{32}\selectfont}

%.................................................................

\tikzposterlatexaffectionproofoff
\usetikzlibrary{shapes.geometric,arrows.meta,positioning}

\usepackage{helvet}
\renewcommand{\familydefault}{\sfdefault}

\definecolor{MyFadedNavy}{rgb}{0.224, 0.259, 0.310}
\definecolor{MyTerracotta}{rgb}{0.914, 0.310, 0.216}
\definecolor{MyIvory}{rgb}{0.965, 0.969, 0.922}

\usetheme{Default}
\definecolorstyle{MyStyle2016}{
    \definecolor{ColorOne}{named}{MyFadedNavy}
    \definecolor{ColorTwo}{named}{MyTerracotta}
    \definecolor{ColorThree}{named}{MyIvory}
}{
    \colorlet{titlebgcolor}{ColorOne}
    \colorlet{titlefgcolor}{white}
    \colorlet{backgroundcolor}{ColorOne!15}
    \colorlet{framecolor}{ColorOne}
    \colorlet{blocktitlebgcolor}{white}
    \colorlet{blocktitlefgcolor}{ColorTwo}
    \colorlet{blockbodybgcolor}{white}
    \colorlet{blockbodyfgcolor}{black}
    \colorlet{innerblocktitlebgcolor}{ColorOne!15}
    \colorlet{innerblocktitlefgcolor}{black}
    \colorlet{innerblockbodybgcolor}{ColorOne!15}
    \colorlet{innerblockbodyfgcolor}{black}
    \colorlet{notebgcolor}{ColorTwo!20}
    \colorlet{notefgcolor}{ColorTwo}
    \colorlet{notefrcolor}{ColorTwo}
}

\usecolorstyle{MyStyle2016}
\useblockstyle[
    titleinnersep=4mm,
    bodyinnersep=4.5mm,
    roundedcorners=12
]{Default}
\useinnerblockstyle[
    titleinnersep=5pt,
    bodyinnersep=7pt
]{Default}

\newcommand{\blocksubheading}[1]{%
    \par\vspace{0.18cm}
    {\color{MyTerracotta}
     \fontsize{31}{34}\selectfont
     \bfseries #1\par}
    \vspace{0.10cm}
}

%.................................................................

\title{\parbox{\linewidth}{\centering
    Addressing a Blind Spot in a
    Topological Data Analysis--Based Data Poisoning Detector}}

\author{\underline{Christian Dane Beels}, Joseph Dorta, Ph.D.}

\institute{Department of Mathematical Sciences,
United States Military Academy, West Point, NY}

\makeatletter
\renewcommand{\TP@maketitle}{%
    \centering\color{titlefgcolor}
    {\fontsize{48}{53}\selectfont
     \bfseries\scshape \@title\par}
    \vspace{0.12cm}
    {\fontsize{28}{33}\selectfont \@author\par}
    \vspace{0.08cm}
    {\fontsize{23}{27}\selectfont \@institute\par}
}
\makeatother

%.................................................................

\begin{document}

\maketitle[
    width=0.97\linewidth,
    innersep=4mm,
    titletotopverticalspace=4mm,
    titletoblockverticalspace=4mm,
    linewidth=0,
    roundedcorners=10
]

\begin{columns}

%=================================================================
% LEFT: visual background and compact controlled design
%=================================================================

\column{0.30}

\block[titleleft,roundedcorners=12]{Motivation}{
    \smallbody\raggedright

    Data poisoning can insert attack-derived traffic with benign
    labels, shifting a learned decision boundary before deployment.
    Monkam et al.\ [1] propose an upstream TDA--OPTICS filter to identify
    poisoned packets before model fitting.

    \vspace{0.08cm}

    {\centering
        \includegraphics[
            width=0.92\linewidth
        ]{Figures/figure_v5_nids_poisoning}\par
    }

    \captionfont\raggedright
    \textbf{Figure 1.} Data poisoning can cause mislabeling, which can
    shift a learned boundary and degrade model performance.
}

\block[titleleft,roundedcorners=12]{Persistent Homology}{
    \smallbody\raggedright

    Persistent homology tracks connected components $(H_0)$ and loops
    $(H_1)$ across filtration scales; long intervals encode durable
    structure.

    \vspace{0.08cm}

    {\centering
        \includegraphics[
            width=0.92\linewidth
        ]{Figures/figure_v1_tda_background}\par
    }

    \captionfont\raggedright
    \textbf{Figure 2.} Point-cloud neighborhoods and the corresponding
    $H_0$/$H_1$ barcode as the scale $\varepsilon$ increases.
}

\block[titleleft,roundedcorners=12]{Reconstruction and Controlled Test}{
    \smallbody\raggedright

    \textbf{Data:} UNSW-NB15 Payload-Byte [2].

    \innerblock{}{
        \centering
        1,500 bytes
        $\to$ 30 $\times$ 50 raster
        $\to \mathrm{Bin}_t$
        $\to$ 2 height $+$ 3 radial filtrations
        $\to$ cubical persistence $(H_0,H_1)$
        $\to \mathbb{R}^{60}$
        $\to$ OPTICS
    }

    \begin{itemize}
        \item Four probe families $\times$ five fixed seeds; each cell has
        5,000 clean $+$ 500 poison packets. Values $=$ means
        $\pm$ population SD.

        \item 4 permutation families: Transpositions (Swap 60 pairs of random bytes), Block Reversal (Mirror 120-byte block), Block Swap (Swap two 60-byte blocks), Cyclic Shift (Rotate every row by 1 index).

        \item Every probe changes the raw vector and preserves its byte multiset.
        \item Ground-truth labels are used for scoring only.
    \end{itemize}
}

%=================================================================
% CENTER: dominant empirical result and compact footer
%=================================================================

\column{0.42}

\block[titleleft,roundedcorners=12]{TDA Pipeline}{
    \centering

    \includegraphics[
        width=0.97\linewidth
    ]{Figures/figure_v9_tda_pipeline}


}

\block[titleleft,roundedcorners=12]{Results}{
    \centering

    \centering
    \innerblock{}{
        \centering
        {\fontsize{34}{39}\selectfont
        \textbf{92–96\%} Fewer Poisoned-Clean Vector Collisions using a multiple-threshold approach.}\\[0.10cm]
        {\fontsize{28}{33}\selectfont}
    }

    \vspace{0.15cm}

    \includegraphics[
        width=0.95\linewidth
    ]{Figures/figure_v6_representation_repair}

    \captionfont\raggedright
    \textbf{Figure 3.} Means over five seeds; 5,000 clean $+$ 500 poison
    packets per cell. Every probe changes the raw byte vector, so each
    collision counted here is one that binarization could in principle
    repair. Nine thresholds cut exact poison--clean matches by 92--96\%.

    {\centering
        
\includegraphics[
            width=0.96\linewidth
        ]{Figures/figure_v7_collision_origin}\par
    }

    \captionfont\raggedright
    \textbf{Figure 4.} Where collisions originate, from a separate audit of
    the single-threshold $\mathbb{R}^{60}$ pipeline. Half the collision mass
    is already present in the raw payload and is unreachable by any
    thresholding scheme; the 46.7\% introduced at binarization is the share
    Figure 3 targets.

}

\block[titleleft,roundedcorners=12]{References and Contact}{
    \begin{minipage}[c]{0.63\linewidth}
        \referencefont\raggedright

        [1] G. F. Monkam, M. J. De Lucia, and N. D. Bastian,
        \emph{Computers \& Security} 144, 103929, 2024.\\[0.04cm]

        [2] N. Moustafa and J. Slay,
        ``UNSW-NB15,'' \emph{MilCIS}, 2015.\\[0.04cm]

        [3] M. Ankerst et al.,
        ``OPTICS,'' \emph{SIGMOD}, pp.\ 49--60, 1999.
    \end{minipage}
    \hfill
    \begin{minipage}[c]{0.34\linewidth}
        \centering

        \begin{minipage}{0.48\linewidth}
            \centering
            
\includegraphics[
                width=2.2in
            ]{Figures/Github_QR.png}\\[-0.03cm]
            \fontsize{15.5}{18}\selectfont
        \end{minipage}
        \hfill
        \begin{minipage}{0.48\linewidth}
            \centering
            
\includegraphics[
                width=2.2in
            ]{Figures/Linkedin_QR.png}\\[-0.03cm]
            \fontsize{15.5}{18}\selectfont
        \end{minipage}
    \end{minipage}
}

%=================================================================
% RIGHT: mathematical mechanism, controlled change, and scope
%=================================================================

\column{0.28}

\block[titleleft,roundedcorners=12]{Why One Cutoff Loses Information}{
    \theorybody\raggedright

    One network packet (Payload Byte) from UNSW-NB15 is represented by $x\in\{0,\ldots,255\}^{n}$, $n=1{,}500$. \\ The maximum byte value
    $M$ is defined as $\max_i x_i = 255$; each byte is then normalized $z=x/M\in[0,1]^{n}$. \\ Define the binarizer at threshold $t$ by
    \[
        \mathrm{Bin}_t(z)_i=
        \begin{cases}
            1, & z_i>t,\\[0.06cm]
            0, & z_i\leq t.
        \end{cases}
    \]
    Let $\sigma\in S_n$ permute byte positions. Since $\sigma$ preserves the
    byte multiset it fixes $M$, so $\sigma z=(\sigma x)/M$ and
    \[
        \mathrm{Bin}_t(\sigma z)=\sigma\,\mathrm{Bin}_t(z).
    \]

    Original pipeline uses a single cutoff, $t=0.4$:
    \[
        A^-=\{i:z_i\leq0.4\},
        \qquad
        A^+=\{i:z_i>0.4\}.
    \]
    Any permutation that acts within these two classes leaves the
    binarization unchanged, so the stabilizer of $z$ is contained in the
    stabilizer of its binarization:
    \[
        G_z:=\operatorname{Stab}(z)
        \;\leq\;
        G_{0.4}:=\operatorname{Stab}\bigl(\mathrm{Bin}_{0.4}(z)\bigr)
        \;=\;S(A^-)\times S(A^+).
    \]
    \begin{center}
        \textbf{Blind spot}: $G_{0.4}\setminus G_z$.
    \end{center}

    Solution: nine thresholds $T=\{0.1,\ldots,0.9\}$. Stack the binarizers
    into a per-position signature
    \[
        Q_T(z)_i=\bigl(\mathrm{Bin}_t(z)_i\bigr)_{t\in T}\in\{0,1\}^{9},
    \]
    and let $B_k$ collect the positions whose signature has exactly $k$
    ones, so $z_i\in\bigl(\tfrac{k}{10},\tfrac{k+1}{10}\bigr]$. The
    stabilizer shrinks to
    \[
        G_z
        \;\leq\;
        G_T:=\operatorname{Stab}\bigl(Q_T(z)\bigr)
        \;=\;\prod_{k=0}^{9}S(B_k)
        \;\leq\;G_{0.4}.
    \]

    \vspace{0.12cm}

    {\color{MyTerracotta}\bfseries
    \begin{center}
        Reducing representation invariance does not guarantee detection, but probability decreases substantially.
    \end{center}\par}
}

\block[titleleft,roundedcorners=12]{Conclusion and Scope}{
    \smallbody\raggedright

    \begin{itemize}
        \item Nine thresholds reduce exact poison--clean collisions by
        92--96\% among raw-changing poison packets, confirming that the
        single-cutoff blind spot is real and closable.

        \item Closing it does \textbf{not} improve detection. The
        540-dimensional representation mainly fragments the OPTICS
        clustering, so poison packets are labeled noise rather than
        recovered as a cluster.

        \item Future work should improve attack separation relative to
        ordinary clean variation before retuning or replacing OPTICS.
    \end{itemize}
}

\end{columns}
\end{document}